\section{Bounded residual forecasting}\label{sec:bounded-decoder}

\subsection{Separate bounded feature blocks}


Each feature block is centered and scaled coordinatewise using A alone;
feature standard deviations are floored at $10^{-6}$. A tilde denotes
this coordinatewise standardization. For a vector $v\in\R^p$, define
\begin{equation}
 T_R(v)=\frac{v}{\sqrt{1+\|v\|_2^2/R^2}},\qquad R=\sqrt p.
 \label{eq:mechanism-main-radial}
\end{equation}
This radial map has norm at most $R$ and is nonexpansive with respect to
its argument. If the common and departure block dimensions are $p_0$
and $p_d$, respectively, set $R_0=\sqrt{p_0}$ and $R_d=\sqrt{p_d}$.
The bounded feature vectors are
\[
 \psi_K(z)=\operatorname{concat}\bigl(
       T_{R_0}(\widetilde\psi_0(z)),T_{R_d}(\widetilde d_K(z))\bigr),
 \qquad
 \psi_R(z)=\operatorname{concat}\bigl(
       T_{R_0}(\widetilde\psi_0(z)),T_{R_d}(\widetilde d_R(z))\bigr).
\]
The two blocks are transformed separately and are not standardized again.
Consequently, appending a departure does not rescale the common features.
Three partitions with 2, 4, and 8 local models, local rank two, and at
most 24 references per local model give $p_0=42$ and $p_d=168$ for the
12-dimensional histories in this experiment. The geometry is fitted
without targets on A and remains fixed during selection and evaluation.

\subsection{Fitting and calibrating the correction}

For horizon $h$, let $f_{0,h}$ be the raw-history ridge baseline. Given
an unstandardized history $x$, let $\psi(x)$ denote the chosen bounded
feature vector after applying its A-fitted input scaler. The forecast is
\begin{equation}
 \widehat y_{t+h}=f_{0,h}(x_t)+\gamma_h q_h(x_t),\qquad
 q_h(x)=b_h+\psi(x)^\top w_h.
 \label{eq:mechanism-main-forecast}
\end{equation}
Here $b_h\in\R$ is a free intercept, $w_h$ is the coefficient vector,
and $\gamma_h$ controls the correction strength. For $n_A$ fitting
histories, the residual targets are
$e_{ih}=y_{i,h}-f_{0,h}(x_i)$, where $y_{i,h}$ is the observed future
at horizon $h$ for fitting history $i$. The correction minimizes
\[
 \frac1{n_A}\sum_{i=1}^{n_A}
       \{e_{ih}-b_h-\psi(x_i)^\top w_h\}^2
       +\alpha_h\|w_h\|_2^2.
\]
Set B selects $\alpha_h\in\{10^{-4},10^{-3},10^{-2},10^{-1},1\}$
by forecast error separately for each horizon. Set C then selects
$\gamma_h\in\{0,1/4,1/2,1\}$. Neither selection step refits the
geometry or decoder, and set D is reserved for final evaluation.

\subsection{A uniform bound including the intercept}

Define the fitting residual mean, centered root mean square, and feature
mean as
\[
 \bar e_h=\frac1{n_A}\sum_i e_{ih},\qquad
 \sigma_{e,h}=\left[\frac1{n_A}\sum_i(e_{ih}-\bar e_h)^2\right]^{1/2},
 \qquad
 \bar\psi=\frac1{n_A}\sum_i\psi(x_i).
\]
\begin{proposition}[Centered ridge correction bound]
\label{prop:mechanism-correction-bound}
For the mean-squared-error objective above with $n_A\geq1$,
$\alpha_h>0$, and an unpenalized intercept, suppose that the feature
map satisfies $\|\psi(x)\|_2\leq R$ on its input domain.
For any scalar $\gamma_h$, the fitted coefficients and correction satisfy
\begin{equation}
 \|w_h\|_2\le\frac{\sigma_{e,h}}{2\sqrt{\alpha_h}},\qquad
 |\gamma_h q_h(x)|\le |\gamma_h|\left[
 |\bar e_h|+(R+\|\bar\psi\|_2)
 \frac{\sigma_{e,h}}{2\sqrt{\alpha_h}}\right].
 \label{eq:mechanism-main-correction-bound}
\end{equation}
\end{proposition}
For a concatenated bounded common and departure vector,
$R=\sqrt{R_0^2+R_d^2}$; for the common-only vector, $R=R_0$.
The free intercept contributes the residual-mean term and cannot be
omitted. This bound controls the correction, not the baseline, total
forecast error, or sensitivity of the preceding reconstruction to changes
in the history.


\paragraph{Proof idea.}
Eliminating the intercept gives
$b_h=\bar e_h-\bar\psi^\top w_h$. If $\Psi_c$ is the matrix whose
rows are $(\psi(x_i)-\bar\psi)^\top$ and $e_c$ is the centered
residual vector, then
$w_h=(\Psi_c^\top\Psi_c+n_A\alpha_h I)^{-1}\Psi_c^\top e_c$.
The coefficient operator has singular values
$s/(s^2+n_A\alpha_h)\leq1/(2\sqrt{n_A\alpha_h})$.
This yields the coefficient bound. Applying Cauchy--Schwarz to
$q_h(x)=\bar e_h+(\psi(x)-\bar\psi)^\top w_h$ yields the
correction bound. Appendix~\ref{app:reconstruction-mechanism} gives
the complete argument, including the nonexpansiveness of the radial map.

The theorem applies separately to each selected horizon and is a property
of bounded-feature ridge regression. A summed-squared-error objective
would need penalty $n_A\alpha_h$ to define the same estimator. The
radial transformation is nonexpansive in its own input, but this does
not give a Lipschitz bound for the preceding signed reconstruction or
for its fallback boundary.

\paragraph{Baseline recovery and predictive gain.}
The calibration grid includes $\gamma_h=0$, which recovers the same
fitted baseline exactly. This is an available forecast rather than a
guarantee that a finite calibration set selects the best future forecast.
The coefficient and correction bounds likewise make no claim about
prediction accuracy. The experiment therefore evaluates error, numerical
coverage, and the correction inequalities as separate outcomes.
